\documentclass[a4paper,12pt]{article}
\usepackage[utf8]{inputenc}
\usepackage[T1]{fontenc}
\usepackage[margin=2.5cm]{geometry}
\usepackage{ctex} % 中文支持
\usepackage{graphicx}
\usepackage{booktabs} % 三线表支持
\usepackage{xcolor}
\usepackage{enumitem} % 列表定制
\usepackage{float}
\usepackage{titlesec}

% -----------------------------------------------------------
% 格式设置区：Level 1 高结构化 + Level 2 标准披露风格
% -----------------------------------------------------------

% 设置层级标题格式（自动编号，移除手动序号）
\titleformat{\section}{\Large\bfseries\color{blue!40!black}}{\thesection}{1em}{}
\titleformat{\subsection}{\large\bfseries\color{gray!80!black}}{\thesubsection}{1em}{}

% 设置段落间距与缩进
\setlength{\parskip}{0.8em}
\setlength{\parindent}{0pt}

% 设置列表间距
\setlist{nosep} % 紧凑列表

% -----------------------------------------------------------
% 文档元数据
% -----------------------------------------------------------
\title{\textbf{2024年度环境、社会及管治（ESG）报告：可持续发展专项章节}}
\author{本公司董事会办公室}
\date{2025年3月}

\begin{document}

\maketitle

% 前言
\section*{前言：以责任驱动高质量发展}

2024年，本公司坚持将可持续发展理念融入企业战略与日常运营。作为国内领先的化工行业上市公司，我们深知合规经营与社会责任不仅是监管的要求，更是企业长青的基石。本章节将对照行业披露标准，用通俗易懂的语言和详实的数据，向各利益相关方汇报我们在环境、社会及管治方面的年度绩效。

% 第一章 环境
\section{环境管理：践行绿色低碳承诺}

我们严格遵守《环境保护法》及相关行业排放标准，致力于降低生产运营对环境的影响。2024年，我们通过加大环保投入和技术升级，在节能减排方面取得了实质性进展。

\subsection{环境绩效概览}

本年度，我们建立了完善的环境数据监测体系。核心环境指标如下表所示：

\begin{table}[H]
\centering
\caption{2024年度环境排放与资源管理绩效表}
\begin{tabular}{@{}ll@{}}
\toprule
\textbf{KPI 指标} & \textbf{2024年绩效数据} \\ \midrule
\textbf{排放管理} & \\
温室气体排放总量 & 39,107.65 吨CO2e \\
排放强度 & 4.58 吨/百万元营收 \\
一般工业固废 & 170.48 吨 \\
危险废物处置率 & 100\% (合规处置) \\
\midrule
\textbf{资源利用} & \\
新鲜水用量 & 149.55 万吨 \\
包装材料使用量 & 7,677 吨 \\
\midrule
\textbf{环保投入} & \\
环保治理总投入 & 457.93 万元 \\
缴纳环保税 & 9.6 万元 \\ \bottomrule
\end{tabular}
\end{table}

\subsection{生产端的清洁化改造}

为了响应国家“双碳”目标，我们在生产源头实施了多项改进措施，重点解决高耗能环节的痛点。

\begin{itemize}
    \item \textbf{热能系统的清洁替代}：在\textbf{核心生产基地}，我们淘汰了传统的燃煤供热方式，完成了清洁能源替代项目。这就好比将厨房里的燃煤灶升级为现代化的清洁能源系统。通过这一举措，该基地实现了生产废气的“零排放”，显著改善了厂区及周边的空气质量。
    \item \textbf{水循环利用体系}：我们旗下化工厂落实了循环经济理念，实施冷凝水回收利用项目。我们将生产蒸汽产生的热净水回收，重新注入锅炉系统。这种“一水多用”的模式，既减少了新鲜水的取用，又节约了再次加热所需的能源。
\end{itemize}

\subsection{物流端的绿色转型}

针对采矿运营中的运输排放问题，我们在\textbf{西北某大型露天煤矿项目}推进了车辆电气化战略，以此降低化石燃料的依赖。

\begin{enumerate}
    \item \textbf{电动化升级}：新增21台纯电动矿卡和36台油电增程式矿卡。这些设备利用电机反转制动技术回收能量，大幅降低了燃油消耗和碳排放。
    \item \textbf{智能化减排}：投入12台无人驾驶电动矿卡。通过精准的电脑控制，车辆运行更平稳，不仅实现了零尾气排放，还大幅降低了作业噪音，减少了对矿区生态的惊扰。
\end{enumerate}

% 第二章 创新
\section{研发创新：赋能产业升级}

科技创新是公司核心竞争力的源泉，也是推动行业安全、高效发展的关键动力。2024年，我们持续加大研发投入，致力于解决行业内的“卡脖子”难题。

\subsection{研发投入与成果}

本年度，我们的研发工作聚焦于本质安全与数智化赋能：
\begin{itemize}
    \item \textbf{资金与人才}：全年研发投入\textbf{4.19亿元}（占营收4.91\%），组建了\textbf{1,064人}的专业研发团队。
    \item \textbf{知识产权}：新增授权专利\textbf{126项}，累计拥有专利596项，通过技术保护构筑了护城河。
\end{itemize}

\subsection{核心技术突破}

我们重点推广了三项具有行业示范意义的技术成果：

\begin{enumerate}
    \item \textbf{现场混装水胶炸药工艺（本质安全）}：
    \begin{itemize}
        \item \textbf{技术特点}：改变了传统炸药成品的运输模式，采用“原材料运输、现场混合”的方式。
        \item \textbf{价值}：此技术填补了国内空白，达到国际先进水平。它不仅消除了运输过程中的爆炸风险，还能根据岩石硬度灵活调节威力，实现了安全与效益的双赢。
    \end{itemize}
    
    \item \textbf{二氧化碳相变膨胀激发管（绿色爆破）}：
    \begin{itemize}
        \item \textbf{技术特点}：利用液态二氧化碳物理膨胀破岩，而非化学燃烧。
        \item \textbf{价值}：国内首创，产能达500万只。由于过程无明火，极大提升了高瓦斯矿井的开采安全性，是典型的“绿色安全技术”。
    \end{itemize}
    
    \item \textbf{智能矿山系统（数智赋能）}：
    \begin{itemize}
        \item \textbf{技术特点}：集成5G、边缘计算与机器人技术。
        \item \textbf{价值}：实现了5G钻机、自动验孔机器人与无人矿卡的协同作业，将工人从高危环境中解放出来，体现了“科技以人为本”的理念。
    \end{itemize}
\end{enumerate}

% 第三章 社会
\section{社会责任：守护员工与社区}

人才是企业最宝贵的资产，社区是企业生存的土壤。我们始终坚持“以人为本，回馈社会”的价值观。

\subsection{职业健康与安全管理}

安全生产是我们对员工最庄严的承诺。2024年，我们共有员工7,399人，劳动合同签订率100\%。我们通过高额投入和严格管理，织密安全防护网。

\begin{table}[H]
\centering
\caption{2024年度安全保障投入表}
\begin{tabular}{@{}ll@{}}
\toprule
\textbf{保障项目} & \textbf{投入金额/措施} \\ \midrule
安全生产责任险 & 355.99 万元 \\
工伤保险 & 892.5 万元 \\
应急管理 & 构建“1+14+N”应急预案体系 \\ \bottomrule
\end{tabular}
\end{table}

\textbf{海外项目最佳实践}：在\textbf{纳米比亚}项目中，面对跨文化管理的挑战，我们建立了标准化的安全管理制度（如“晨会制度”和“手指口述”法）。这一严谨的管理体系帮助项目部实现了\textbf{连续130万小时无损工事件}的优异成绩，证明了中国企业的安全管理水平已达到国际标准。

\subsection{企业公民责任}

作为央企控股企业，我们积极响应国家乡村振兴号召，并在突发灾害面前勇于担当。

\begin{itemize}
    \item \textbf{公益捐赠}：全年对外捐赠及公益总投入\textbf{1,219万元}，其中1,125万元专项用于乡村振兴，助力改善农村基础设施和民生。
    \item \textbf{抢险救灾案例}：2024年\textbf{湖南某地}发生决堤险情时，我们迅速启动应急响应机制。依托专业爆破技术，我们采用“松动爆破”方案，在极短时间内开采\textbf{8.4万吨}石料支援封堵行动。这次行动不仅是物资的支援，更是专业技术在社会急难险重任务中的成功应用。
\end{itemize}

% 第四章 治理
\section{公司治理：合规与价值创造}

完善的公司治理是企业实现可持续发展的制度保障。作为\textbf{上级集团}成员企业，我们致力于维护投资者权益，打造廉洁透明的供应链。

\subsection{治理架构与投资者回报}

我们构建了科学的决策机制和积极的分红政策：
\begin{itemize}
    \item \textbf{董事会结构}：采用“3+3+3”模式（3名股东董事、3名独立董事、3名执行董事），确保了决策的独立性与制衡性。公司因此获得交易所信披考评“A级”。
    \item \textbf{股东回报}：2024年，公司共派发现金红利\textbf{2.54亿元}（每10股2.05元），分红比例达可分配利润的\textbf{40.12\%}，切实让投资者分享企业发展红利。
\end{itemize}

\subsection{供应链廉洁合规}

我们高度重视供应链的商业道德风险管理。针对合作的1,548家供应商，我们推行了强制性的合规准入机制：
\begin{itemize}
    \item \textbf{廉洁协议}：要求所有供应商签署《廉洁协议》，签署率达到\textbf{100\%}。
    \item \textbf{合规管控}：通过制度约束，确保采购过程公开、透明，杜绝商业贿赂，共同营造健康的行业生态。
\end{itemize}

\end{document}